\section{The explicit isomorphism}
\label{sec:explicit-isomorphism}

Section~\ref{sec:fiber-product} identifies the abstract group suggested
by the native invariants. We now strengthen that classification by
constructing an explicit isomorphism
\[
\Phi\colon
\operatorname{Aut}(G60)
\longrightarrow
S_5\times_{C_2}D_8.
\]

This step removes any remaining dependence on fingerprint comparison.
The two groups are connected element by element, and multiplication is
verified throughout the complete multiplication table.

\subsection{Native and model groups}

Let
\[
A=\operatorname{Aut}(G60)
\]
be represented by its $480$ native permutations on the fixed
sixty-vertex labeling.

Let
\[
F=S_5\times_{C_2}D_8
\]
be the fiber-product model of
Theorem~\ref{thm:fiber-product-identification}. Its elements are pairs
\[
(\sigma,d)\in S_5\times D_8
\]
satisfying
\[
\operatorname{sgn}(\sigma)=\chi_{\mathrm{V4}}(d).
\]

Both groups are enumerated exactly and use deterministic element
encodings. The native elements are stored as permutation tuples, while
the model elements are stored as compatible coordinate pairs.

\subsection{Generator correspondence}

A generating set
\[
\mathcal{S}_A=\{g_1,\ldots,g_k\}
\]
is selected from the native permutation group. A corresponding set
\[
\mathcal{S}_F=\{h_1,\ldots,h_k\}
\]
is selected in the fiber-product model.

The correspondence
\[
g_i\longmapsto h_i
\]
is chosen so that the native generators and model generators satisfy the
same multiplication relations.

The generator images include representatives of the following
structural features:

\begin{itemize}
\item the normal $A_5$ subgroup;
\item odd $S_5$ parity;
\item the central involution;
\item an order-four dihedral direction;
\item a reflection direction.
\end{itemize}

The exact generator arrays and model coordinates are recorded in the
certificate described in Appendix~B.

\subsection{Extension to all elements}

Starting from the identity, the generator correspondence is extended
through simultaneous traversal of the two Cayley graphs.

Suppose that a native element $x\in A$ has already been assigned an
image $\Phi(x)\in F$. For each paired generator
\[
g_i\longmapsto h_i,
\]
define
\[
\Phi(xg_i)=\Phi(x)h_i.
\]

Whenever an element is reached by more than one word, the independently
computed model images are required to agree. Thus the traversal checks
the consistency of the generator correspondence while constructing the
map.

The traversal reaches every element of $A$ and every element of $F$.
It produces a bijection
\[
\Phi\colon A\longrightarrow F
\]
with
\[
|\operatorname{dom}\Phi|
=
|\operatorname{im}\Phi|
=
480.
\]

No native element is assigned two different model images, and no two
native elements are assigned the same model image.

\subsection{Complete multiplication verification}

After constructing the bijection, we test the homomorphism law
independently for every ordered pair
\[
(x,y)\in A\times A.
\]

For each pair, the verification checks
\[
\Phi(xy)=\Phi(x)\Phi(y).
\]

The number of ordered pairs is
\[
|A|^2
=
480^2
=
230400.
\]

All $230400$ checks pass.

\begin{theorem}
\label{thm:explicit-isomorphism}
The exported map
\[
\Phi\colon
\operatorname{Aut}(G60)
\longrightarrow
S_5\times_{C_2}D_8
\]
is a group isomorphism.
\end{theorem}

\begin{proof}
The simultaneous generator traversal constructs a bijection between the
two $480$-element sets. The exhaustive multiplication test verifies
\[
\Phi(xy)=\Phi(x)\Phi(y)
\]
for every ordered pair $(x,y)\in A\times A$. Therefore $\Phi$ is a
bijective homomorphism and hence an isomorphism.
\end{proof}

\subsection{Certificate digest}

The full mapping is stored as a machine-readable artifact. A canonical
serialization of the mapping has SHA-256 digest
\[
\texttt{9a2b9bc2294a858cb7411fec33aa9c69786ad387184de839925f73453a3e169d}.
\]

The digest binds the theorem statement to one exact exported
correspondence. It does not replace the mathematical proof, but it
allows the explicit certificate used in the verification to be
identified without ambiguity.

\subsection{Consequences}

The explicit isomorphism simultaneously confirms:

\begin{enumerate}
\item the native automorphism group has no missing elements;
\item the fiber-product model has no extra elements;
\item the proposed generator correspondence is globally consistent;
\item the group law is preserved on all elements, not merely on the
generators or on selected relations.
\end{enumerate}

The classification
\[
\operatorname{Aut}(G60)
\cong
S_5\times_{C_2}D_8
\]
is therefore established both structurally and constructively.
